\section*{Appendix}
\label{sec:appendix}

\subsection*{A. Additional Experimental Details}

\textbf{Feature caching.} Features were extracted once per backbone on the full Waterbirds and CelebA splits, stored as \texttt{.pt} tensors, and loaded for all probe experiments. This ensures identical features across all seeds and avoids any stochasticity from the extraction step.

\textbf{Group-balanced sampling.} For Waterbirds, the four groups are \{landbird on land, landbird on water, waterbird on land, waterbird on water\}. The balanced validation probe training set samples equal numbers from each group. For CelebA, the four groups are \{Blond Female, Non-Blond Female, Blond Male, Non-Blond Male\}.

\textbf{Reproducibility.} All probes were run with 5 random seeds (0--4) for the logistic regression solver initialization. ANOVA and pairwise $t$-tests are performed across seed-level ratio values ($n=5$ per paradigm per dataset).

\subsection*{B. Background-Replacement Augmentation Details}

The SimCLR-NoBackground condition replaces background pixels (identified via CUB-200-2011 segmentation masks) with patches sampled from Places365 images before applying standard SimCLR augmentations (random crop, horizontal flip, color jitter, grayscale). The pixel difference metric is computed as the mean absolute per-pixel difference between the NoBackground and Original augmented view pairs, averaged over a held-out mini-batch of 256 images. The reported value of 0.9656 vastly exceeds the 0.05 activation threshold, confirming the mechanism is active.
